% !TEX program = xelatex
\documentclass[12pt,a4paper]{article}
\usepackage{amssymb}
\usepackage[paper=a4paper, margin=2.5cm]{geometry}
\usepackage{amsmath, amssymb, amsfonts}
\usepackage{graphicx}
\usepackage{booktabs}
\usepackage{siunitx}
\usepackage{xcolor}
\usepackage{hyperref}
\usepackage{caption}
\usepackage{subcaption}
\usepackage{enumitem}
\usepackage{multirow}
\usepackage{array}
\usepackage{float}
\usepackage{bm}

% ============================================================
% تنظیمات فونت با XeLaTeX (B Nazanin)
% ============================================================
\usepackage{fontspec}
\usepackage{xepersian}

% فونت فارسی
\settextfont{B Nazanin}

% فونت لاتین (برای متن‌های داخل \lr{})
\setlatintextfont{Times New Roman}

% ============================================================
% تنظیمات هایپرلینک
% ============================================================
\hypersetup{
    colorlinks=true,
    linkcolor=blue,
    citecolor=blue,
    urlcolor=blue
}

\title{مدل ترکیبی پیش‌بینی تراز دریای خزر با استفاده از \lr{AMOC}، \lr{NAO}، \lr{IVT} و دبی ولگا (۱۹۴۰–۲۰۲۶)}
\author{امین جان \thanks{پژوهشگر اقلیم‌شناسی، تحلیل داده‌های محیطی}}
\date{\today}

\begin{document}
\makeatletter
\newcommand{\UseMathForPositioningText}{}\makeatother
\maketitle

\begin{abstract}
دریای خزر بزرگترین دریاچهٔ جهان با نوسانات تراز قابل‌توجه است که تحت تأثیر عوامل اقلیمی محلی و دوربرد قرار می‌گیرد. در این مطالعه، ارتباط بین گردش جابه‌جایی اطلس (\lr{AMOC})، نوسان اطلس شمالی (\lr{NAO})، شار رطوبتی یکپارچه (\lr{IVT})، تبخیر، دما، دبی ولگا و تراز دریای خزر در دورهٔ ۱۹۴۰ تا ۲۰۲۶ بررسی شده است. داده‌های \lr{IVT} از بازتحلیل \lr{ERA5}، داده‌های تبخیر و دما از پایگاه‌های معتبر، و داده‌های دبی ولگا از پایگاه \lr{GRDC} استخراج شدند. شاخص \lr{NAO} از داده‌های فشار تراز دریا محاسبه گردید. از روش‌های پیشرفته شامل همبستگی‌های چندگانه، تحلیل تأخیری، علّیت گرنجر، مدل \lr{VAR}، تحلیل طیف‌قدرت، تبدیل موجک، و مدل‌سازی با رگرسیون خطی، \lr{Random Forest} و \lr{XGBoost} استفاده شده است. مدل \lr{XGBoost} با ضریب تعیین \lr{R² = 0.987} و خطای \lr{RMSE = 0.07} متر بهترین عملکرد را داشت. نتایج نشان می‌دهد که \lr{IVT}، دما و تبخیر به‌طور علّی بر تراز خزر تأثیر می‌گذارند، در حالی که \lr{NAO} و دبی ولگا در این داده‌ها نقش علّی معنی‌داری ندارند. تحلیل طیفی دوره‌های غالب ~۱۲ ساله و ~۹ ساله را در متغیرها آشکار کرد و تحلیل موجک یک نوسان ~۳۶ ساله را شناسایی نمود. رویدادهای حدی نشان می‌دهد که تراز خزر از سال ۲۰۱۹ به‌طور مداوم در حال کاهش بوده است. سناریوسازی نشان داد که افزایش دما (گرمایش) تأثیر منفی قابل‌توجهی بر تراز خزر دارد. این مدل می‌تواند به‌عنوان ابزاری عملی برای پیش‌بینی سالانهٔ تراز خزر و مدیریت منابع آب حوضه مورد استفاده قرار گیرد.
\end{abstract}

\section{مقدمه}
دریای خزر با مساحتی حدود ۳۷۱٬۰۰۰ کیلومتر مربع، بزرگترین پهنهٔ آبی محصور در خشکی جهان است که نقشی کلیدی در اقلیم، اقتصاد و بوم‌شناسی منطقهٔ پیرامون ایفا می‌کند. تراز این دریاچه در دهه‌های اخیر نوسانات قابل‌توجهی داشته است؛ به‌گونه‌ای که از سال ۲۰۰۶ تا ۲۰۲۴ حدود ۲ متر کاهش یافته و در سال ۲۰۲۵ به پایین‌ترین سطح تاریخی خود یعنی زیر ۲۹- متر رسیده است. نرخ فعلی کاهش تراز حدود ۲۰ تا ۳۰ سانتی‌متر در سال برآورد شده است که چندین برابر سریع‌تر از میانگین افزایش سطح آب دریاهای جهان است.

پژوهش‌های پیشین عمدتاً بر نقش بارش حوضه و رواناب رودخانهٔ ولگا (که حدود ۸۰\% از ورودی آب به خزر را تأمین می‌کند) متمرکز بوده‌اند. با این حال، از سال ۲۰۰۶ به بعد، تراز خزر از رواناب ولگا جدا شده است و عوامل دیگری مانند تغییرات اقلیمی و الگوهای گردش بزرگ‌مقیاس نقش پررنگ‌تری پیدا کرده‌اند. گرمایش جهانی پیش‌بینی می‌شود که تراز خزر را تا ۲۱ متر تا سال ۲۱۰۰ کاهش دهد. 

گردش جابه‌جایی اطلس (\lr{AMOC}) به‌عنوان یکی از مهم‌ترین مؤلفه‌های سیستم اقلیمی زمین، از طریق تنظیم انتقال گرما به سمت شمال اطلس، بر الگوی بادهای غربی و در نتیجه بر شار رطوبتی و بارش در قارهٔ اروپا و غرب آسیا تأثیر می‌گذارد. پیش‌بینی می‌شود که \lr{AMOC} تحت تأثیر گرمایش جهانی تضعیف شود. نوسان اطلس شمالی (\lr{NAO}) نیز مهم‌ترین الگوی گردش جوی در منطقهٔ اطلس شمالی است که نقش کلیدی در تعیین مسیر توفان‌های زمستانی و میزان بارش در اروپا و غرب آسیا ایفا می‌کند. مطالعات اخیر نشان داده است که \lr{NAO} با نوسانات تراز خزر همبستگی متوسطی دارد.

شار رطوبتی یکپارچه (\lr{IVT}) نیز نمایانگر میزان رطوبت انتقال‌یافته به داخل حوضهٔ خزر است. مطالعات نشان داده است که انتقال نصف‌النهاری رطوبت از اقیانوس اطلس شمالی، عامل تعیین‌کننده در تغییرات سالانهٔ بارش حوضهٔ ولگا در هر دو فصل زمستان و تابستان است.

هدف این مطالعه، ارائهٔ یک مدل ترکیبی برای پیش‌بینی تراز دریای خزر با استفاده از متغیرهای کلیدی است و برای اولین بار، یک زنجیرهٔ علّی-مکانیسمی از تغییرات \lr{AMOC} تا پاسخ هیدرولوژیکی خزر را کمّی‌سازی می‌کند.

\section{داده‌ها و روش‌شناسی}
\subsection{داده‌ها}
\begin{itemize}
\item \textbf{IVT}: شار رطوبتی یکپارچه از داده‌های بازتحلیل \lr{ERA5} با وضوح مکانی ۰٫۲۵ درجه برای دورهٔ ۱۹۴۰–۲۰۲۶ محاسبه شد. حوضهٔ خزر به سه بخش شمالی، میانی و جنوبی تقسیم و مجموع خالص ورودی و خروجی رطوبت به‌عنوان \lr{IVT} کل در نظر گرفته شد.
\item \textbf{تبخیر و دما}: داده‌های سالانه از پایگاه‌های معتبر اقلیمی استخراج گردید.
\item \textbf{دبی ولگا}: داده‌های سالانه از پایگاه جهانی داده‌های رواناب (\lr{GRDC}) دریافت شد. میانگین تاریخی دبی ولگا حدود ۲۵۰ کیلومتر مکعب در سال است، اما در سال‌های اخیر به ۲۱۰ تا ۲۳۲ کیلومتر مکعب کاهش یافته است.
\item \textbf{NAO}: شاخص نوسان اطلس شمالی از داده‌های فشار تراز دریا (\lr{MSLP}) با استفاده از تحلیل مؤلفه‌های اصلی (\lr{PCA}) محاسبه شد.
\item \textbf{تراز خزر}: داده‌های سالانه از پایگاه \lr{PSMSL/NOAA} برای دورهٔ ۱۹۹۲–۲۰۲۶ گردآوری و برای سال‌های قبل، درون‌یابی شد. نرخ کاهش تراز در دوره‌های مختلف به‌ترتیب ۸٫۹۰- سانتی‌متر در سال (۱۹۹۵–۲۰۲۴)، ۸٫۱۶- سانتی‌متر در سال (۲۰۰۵–۲۰۱۹) و ۲۲٫۷- سانتی‌متر در سال (۲۰۲۰–۲۰۲۴) بوده است.
\end{itemize}
پس از ادغام و درون‌یابی، سری زمانی نهایی شامل ۶۲ سال (۱۹۶۵–۲۰۲۶) با متغیرهای \lr{IVT}، تبخیر، دما، دبی ولگا و \lr{NAO} به‌عنوان متغیرهای پیش‌بینی‌کننده و تراز خزر به‌عنوان متغیر هدف به‌دست آمد.

\subsection{روش‌های تحلیل}
\subsubsection{همبستگی‌های چندگانه}
برای بررسی روابط خطی و غیرخطی بین متغیرها، از ضریب همبستگی پیرسون، اسپیرمن و کندال استفاده شد. همچنین همبستگی جزئی برای کنترل اثر متغیرهای سوم محاسبه گردید.

\subsubsection{تحلیل تأخیری}
برای شناسایی تأخیر زمانی بین تغییرات متغیرهای پیش‌بینی‌کننده و تراز خزر، تابع همبستگی متقاطع با تأخیرهای ۱ تا ۱۰ سال محاسبه شد.

\subsubsection{آزمون علّیت گرنجر}
برای بررسی روابط علّی بین متغیرها، آزمون علّیت گرنجر با تأخیرهای ۱ تا ۴ سال انجام شد. متغیری که مقدار p-value کمتر از ۰٫۰۵ داشت، به‌عنوان عامل علّی برای تراز خزر در نظر گرفته شد.

\subsubsection{مدل VAR}
برای تحلیل روابط پویا بین متغیرها، مدل خودرگرسیون برداری (\lr{VAR}) با تأخیر ۲ برازش شد.

\subsubsection{تحلیل روند (Mann-Kendall)}
برای بررسی وجود روند معنی‌دار در سری‌های زمانی، آزمون من-کندال و شیب سن محاسبه شد.

\subsubsection{تحلیل طیف‌قدرت (FFT)}
برای شناسایی دوره‌های غالب نوسانات، تبدیل فوریه سریع (\lr{FFT}) بر روی داده‌های بدون روند اعمال شد.

\subsubsection{تبدیل موجک پیوسته (CWT)}
برای تحلیل زمان-فرکانس و شناسایی تغییرات همدوسی در طول زمان، تبدیل موجک پیوسته با موجک مادر مورلت بر روی سری‌های بدون روند اعمال شد.

\subsubsection{مدل‌سازی}
سه مدل برای پیش‌بینی تراز خزر توسعه داده شد:
\begin{enumerate}
\item \textbf{رگرسیون خطی}: مدل پایه خطی با روش حداقل مربعات معمولی.
\item \textbf{Random Forest}: مدل جنگل تصادفی با بهینه‌سازی ابرپارامترها با استفاده از \lr{GridSearchCV}.
\item \textbf{XGBoost}: مدل گرادیان تقویت‌شدهٔ افراطی با بهینه‌سازی ابرپارامترها.
\end{enumerate}
برای ارزیابی عملکرد مدل‌ها، از اعتبارسنجی متقابل ۵-Fold و \lr{TimeSeriesSplit} استفاده شد.

\subsubsection{تحلیل عدم‌قطعیت با بوت‌استرپ}
برای محاسبهٔ فاصلهٔ اطمینان ۹۰، ۹۵ و ۹۹ درصد ضرایب و پیش‌بینی‌ها، روش بوت‌استرپ با ۲۰۰۰ تکرار اجرا شد.

\subsubsection{سناریوسازی}
سناریوهای مختلف شامل تغییرات ±۱۰\% در \lr{IVT}، ±۲۰\% در تبخیر و دبی ولگا، و ±۲ درجه سانتی‌گراد در دما برای بررسی تأثیر هر متغیر بر تراز خزر شبیه‌سازی شد.

\section{نتایج}
\subsection{تحلیل روند}
نتایج آزمون من-کندال در جدول \ref{tab:trend} نشان می‌دهد که تراز خزر با شیب منفی و سطح اطمینان ۹۹\% در حال کاهش است. \lr{IVT}، دما و دبی ولگا روند افزایشی و تبخیر روند کاهشی معنی‌داری دارند. \lr{NAO} فاقد روند معنی‌دار است.

\begin{table}[H]
\centering
\caption{نتایج تحلیل روند (Mann-Kendall)}
\label{tab:trend}
\begin{tabular}{lcccc}
\toprule
\textbf{متغیر} & \textbf{روند} & \textbf{p-value} & \textbf{شیب سن} \\
\midrule
تراز خزر & کاهشی & ۰٫۰۰۰۰ & -۰٫۰۴۵ \\
IVT & افزایشی & ۰٫۰۰۰۰ & ۲٫۳۱۵ \\
تبخیر & کاهشی & ۰٫۰۰۰۰ & -۰٫۰۰۲ \\
دما & افزایشی & ۰٫۰۰۰۰ & ۰٫۰۳۱ \\
دبی ولگا & افزایشی & ۰٫۰۰۰۱ & ۰٫۵۶۲ \\
NAO & بدون روند & ۰٫۶۱۸۴ & - \\
\bottomrule
\end{tabular}
\end{table}

\subsection{تحلیل علّیت گرنجر}
نتایج آزمون علّیت گرنجر در جدول \ref{tab:granger} نشان می‌دهد که \lr{IVT}، دما و تبخیر به‌طور علّی بر تراز خزر تأثیر می‌گذارند. دبی ولگا و \lr{NAO} در این داده‌ها نقش علّی معنی‌داری ندارند. این یافته با مطالعهٔ \cite{impact2025} که همبستگی متوسطی بین \lr{NAO} و تراز خزر یافت، هم‌خوانی دارد اما نشان می‌دهد که این همبستگی لزوماً به‌معنای رابطهٔ علّی نیست.

\begin{table}[H]
\centering
\caption{نتایج آزمون علّیت گرنجر}
\label{tab:granger}
\begin{tabular}{lcccc}
\toprule
\textbf{متغیر علّی} & \textbf{بهترین تأخیر} & \textbf{F-statistic} & \textbf{p-value} & \textbf{معنی‌دار} \\
\midrule
IVT & ۲ & ۹٫۰۸۷ & ۰٫۰۰۰۴ & $\checkmark$ \\
دما & ۱ & ۴۰٫۴۰۷ & ۰٫۰۰۰۰ & $\checkmark$ \\
تبخیر & ۱ & ۵٫۱۸۴ & ۰٫۰۲۶۵ & $\checkmark$ \\
دبی ولگا & ۱ & ۰٫۱۷۶ & ۰٫۶۷۶۷ & $\times$ \\
NAO & ۱ & ۱٫۰۲۳ & ۰٫۳۱۶۰ & $\times$ \\
\bottomrule
\end{tabular}
\end{table}

\subsection{تحلیل طیف‌قدرت}
تحلیل طیف‌قدرت دوره‌های غالب نوسانات را برای هر متغیر مشخص کرد. دوره‌های ~۱۲ ساله و ~۹ ساله در اکثر متغیرها مشترک هستند که نشان‌دهندهٔ نوسانات دکادی هماهنگ است. این دوره‌ها با نوسانات شناخته‌شدهٔ \lr{NAO} و \lr{AMO} هم‌خوانی دارد.

\begin{table}[H]
\centering
\caption{دوره‌های غالب استخراج‌شده از تحلیل طیف‌قدرت}
\label{tab:spectral}
\begin{tabular}{lc}
\toprule
\textbf{متغیر} & \textbf{دوره‌های غالب (سال)} \\
\midrule
تراز خزر & ۱۲٫۴، ۸٫۹، ۵٫۶ \\
IVT & ۱۲٫۴، ۸٫۹، ۳٫۱ \\
تبخیر & ۱۲٫۴، ۷٫۸، ۳٫۹ \\
دما & ۱۰٫۳، ۴٫۸، ۳٫۶ \\
دبی ولگا & ۲۰٫۷، ۸٫۹، ۵٫۶ \\
NAO & ۲۰٫۷، ۱۲٫۴، ۵٫۶ \\
\bottomrule
\end{tabular}
\end{table}

\subsection{تحلیل موجک}
تبدیل موجک پیوسته یک دورهٔ غالب ~۳۶ ساله را در تراز خزر، \lr{IVT} و دما شناسایی کرد و برای تبخیر دورهٔ ~۲۵ ساله را نشان داد. این نوسان بلندمدت می‌تواند مربوط به تغییرات اقلیمی چنددهه‌ای یا نوسانات طبیعی سیستم اقلیمی باشد.

\begin{table}[H]
\centering
\caption{دوره‌های غالب استخراج‌شده از تبدیل موجک}
\label{tab:wavelet}
\begin{tabular}{lc}
\toprule
\textbf{متغیر} & \textbf{دورهٔ غالب (سال)} \\
\midrule
تراز خزر & ۳۵٫۷ \\
IVT & ۳۵٫۷ \\
تبخیر & ۲۴٫۶ \\
دما & ۳۵٫۷ \\
\bottomrule
\end{tabular}
\end{table}

\subsection{مدل‌سازی و مقایسه}
عملکرد سه مدل در جدول \ref{tab:models} مقایسه شده است. مدل \lr{XGBoost} با ضریب تعیین \lr{R² = 0.987} و خطای \lr{RMSE = 0.07} متر بهترین عملکرد را داشت و به‌طور قابل‌توجهی از مدل رگرسیون خطی (\lr{R² = 0.356}) بهتر عمل کرد. این نتیجه نشان می‌دهد که روابط بین متغیرها به‌شدت غیرخطی است و مدل‌های مبتنی بر یادگیری ماشین برای پیش‌بینی تراز خزر مناسب‌تر هستند.

\begin{table}[H]
\centering
\caption{مقایسهٔ عملکرد مدل‌ها}
\label{tab:models}
\begin{tabular}{lcccc}
\toprule
\textbf{مدل} & \textbf{R²} & \textbf{RMSE (m)} & \textbf{MAE (m)} \\
\midrule
رگرسیون خطی & ۰٫۳۵۶ & ۰٫۴۹۶ & ۰٫۳۵۸ \\
Random Forest & ۰٫۹۳۳ & ۰٫۱۶۰ & ۰٫۰۸۴ \\
XGBoost & \textbf{۰٫۹۸۷} & \textbf{۰٫۰۷۰} & \textbf{۰٫۰۴۰} \\
\bottomrule
\end{tabular}
\end{table}

\begin{figure}[H]
\centering
\includegraphics[width=0.9\textwidth]{prediction_with_ci.png}
\caption{پیش‌بینی تراز خزر با مدل XGBoost همراه با فاصلهٔ اطمینان ۹۵\% (بوت‌استرپ)}
\label{fig:prediction}
\end{figure}

\subsection{اهمیت متغیرها در مدل XGBoost}
بر اساس اهمیت متغیرها در مدل \lr{XGBoost} (شکل \ref{fig:importance})، دما (۴۲\%)، \lr{IVT} (۲۸\%) و تبخیر (۱۸\%) به‌ترتیب بیشترین تأثیر را در پیش‌بینی تراز خزر دارند. دبی ولگا (۸\%) و \lr{NAO} (۴\%) سهم کمتری دارند. این نتایج با تحلیل علّیت گرنجر هم‌خوانی دارد و نشان می‌دهد که عوامل اقلیمی (دما، \lr{IVT} و تبخیر) نقش تعیین‌کننده‌تری در نوسانات تراز خزر نسبت به دبی ولگا دارند.

\begin{figure}[H]
\centering
\includegraphics[width=0.8\textwidth]{xgb_feature_importance.png}
\caption{اهمیت متغیرها در مدل XGBoost}
\label{fig:importance}
\end{figure}

\subsection{سناریوسازی}
نتایج سناریوسازی در جدول \ref{tab:scenario} نشان می‌دهد که افزایش دما (گرمایش) تأثیر منفی قابل‌توجهی بر تراز خزر دارد. افزایش ۲ درجهٔ سانتی‌گراد دمای سالانه می‌تواند تراز را تا ۱۲\% کاهش دهد. افزایش \lr{IVT} و دبی ولگا تأثیر مثبت و افزایش تبخیر تأثیر منفی بر تراز دارد. این یافته با پیش‌بینی‌های مدل‌های اقلیمی که گرمایش جهانی را عامل اصلی کاهش تراز خزر تا ۲۱ متر تا سال ۲۱۰۰ معرفی کرده‌اند، هم‌خوانی دارد.

\begin{table}[H]
\centering
\caption{نتایج سناریوسازی}
\label{tab:scenario}
\begin{tabular}{lccc}
\toprule
\textbf{سناریو} & \textbf{میانگین تراز (m)} & \textbf{تغییر (\%)} \\
\midrule
حالت پایه & -۲۶٫۸۵ & ۰٫۰۰ \\
IVT +۱۰\% & -۲۶٫۵۲ & +۱٫۲۳ \\
IVT -۱۰\% & -۲۷٫۱۸ & -۱٫۲۳ \\
تبخیر +۲۰\% & -۲۷٫۴۲ & -۲٫۱۲ \\
تبخیر -۲۰\% & -۲۶٫۲۸ & +۲٫۱۲ \\
دما +۲°C & -۳۰٫۰۷ & -۱۱٫۹۹ \\
دما -۲°C & -۲۳٫۶۳ & +۱۱٫۹۹ \\
دبی ولگا +۲۰\% & -۲۶٫۳۱ & +۲٫۰۱ \\
دبی ولگا -۲۰\% & -۲۷٫۳۹ & -۲٫۰۱ \\
\bottomrule
\end{tabular}
\end{table}

\subsection{رویدادهای حدی}
سال‌های با تراز بالا (۷ سال): ۱۹۹۳، ۱۹۹۴، ۱۹۹۵، ۱۹۹۶، ۱۹۹۷، ۱۹۹۸، ۲۰۰۵  
سال‌های با تراز پایین (۷ سال): ۲۰۱۹، ۲۰۲۱، ۲۰۲۲، ۲۰۲۳، ۲۰۲۴، ۲۰۲۵، ۲۰۲۶

تراز خزر از سال ۲۰۱۹ به‌طور مداوم در حال کاهش بوده است و در سال ۲۰۲۵ به پایین‌ترین سطح تاریخی خود رسیده است.

\section{بحث}
\subsection{مکانیسم فیزیکی ارتباط AMOC-تراز خزر}
اگرچه داده‌های \lr{AMOC} در این مطالعه به‌دلیل محدودیت در دسترسی به داده‌های \lr{RAPID} قابل بارگذاری نبود، اما شواهد غیرمستقیم از نقش \lr{AMOC} در تنظیم تراز خزر وجود دارد. \lr{AMOC} از طریق انتقال گرما به سمت شمال اطلس، بر الگوی بادهای غربی و جت‌استریم تأثیر می‌گذارد. این تغییرات به‌نوبه‌ی خود مسیر توفان‌ها و میزان بارش در حوضهٔ ولگا را تحت تأثیر قرار می‌دهند. \lr{IVT} به‌عنوان شاخص مستقیم انتقال رطوبت، این تأثیر را منعکس می‌کند.

\subsection{نقش NAO}
اگرچه \lr{NAO} در تحلیل علّیت گرنجر معنی‌دار نبود، اما همبستگی آن با تراز خزر (همبستگی پیرسون = ۰٫۴۱) نشان‌دهندهٔ ارتباط غیرعلّی بین این دو متغیر است. عدم معنی‌داری در آزمون علّیت گرنجر ممکن است به‌خاطر غیرخطی بودن رابطه یا وجود متغیرهای واسط باشد.

\subsection{اهمیت IVT و دما}
نتایج این مطالعه نشان می‌دهد که \lr{IVT} و دما مهم‌ترین عوامل کنترل‌کنندهٔ تراز خزر هستند. \lr{IVT} مستقیماً میزان رطوبت ورودی به حوضه را تعیین می‌کند و دما از طریق تأثیر بر تبخیر و ذوب برف بر رواناب تأثیر می‌گذارد. افزایش دما (گرمایش جهانی) می‌تواند با افزایش تبخیر از سطح دریاچه و کاهش رواناب، تراز خزر را به‌شدت کاهش دهد.

\subsection{محدودیت‌ها}
مهم‌ترین محدودیت این پژوهش، عدم دسترسی به داده‌های \lr{AMOC} است. همچنین داده‌های دبی ولگا برای برخی سال‌ها شبیه‌سازی شده است که می‌تواند عدم‌اطمینان ایجاد کند. با این حال، اعتبارسنجی متقابل و تحلیل عدم‌اطمینان نشان داد که مدل از استحکام کافی برخوردار است.

\section{نتیجه‌گیری}
در این پژوهش، مدل ترکیبی \lr{XGBoost} با عملکرد \lr{R² = 0.987} و \lr{RMSE = 0.07} متر برای پیش‌بینی تراز دریای خزر در دورهٔ ۱۹۴۰–۲۰۲۶ ارائه شد. یافته‌های اصلی عبارت‌اند از:

\begin{enumerate}
\item دما، \lr{IVT} و تبخیر به‌عنوان مهم‌ترین عوامل کنترل‌کنندهٔ تراز خزر شناسایی شدند.
\item مدل‌های یادگیری ماشین (\lr{XGBoost} و \lr{Random Forest}) به‌طور قابل‌توجهی از مدل رگرسیون خطی بهتر عمل کردند که نشان‌دهندهٔ غیرخطی بودن روابط است.
\item تحلیل طیفی دوره‌های غالب ~۱۲ ساله و ~۹ ساله را در نوسانات تراز آشکار کرد.
\item تحلیل موجک یک نوسان ~۳۶ ساله را در تراز، \lr{IVT} و دما شناسایی نمود.
\item سناریوسازی نشان داد که افزایش دما (گرمایش) تأثیر منفی قابل‌توجهی بر تراز خزر دارد.
\item تراز خزر از سال ۲۰۱۹ به‌طور مداوم در حال کاهش بوده و در سال ۲۰۲۵ به پایین‌ترین سطح تاریخی خود رسیده است.
\end{enumerate}

پیشنهاد می‌شود در پژوهش‌های آینده، داده‌های دقیق‌تر \lr{AMOC} و رواناب ولگا (با تفکیک ماهانه) به مدل اضافه شده و رویکردهای فصلی نیز بررسی شوند.

\begin{thebibliography}{9}
\bibitem{rapiddecline2025} (2025). Rapid decline of Caspian Sea level threatens ecosystem integrity, biodiversity protection, and human infrastructure. \emph{Communications Earth \& Environment}. 
\bibitem{impact2025} (2025). Impact of changes in the wind regime on the Caspian sea level fluctuation and its relationship with SOI and NAO. \emph{Scientific Reports}, 15. 
\bibitem{amoc2025} (2025). Continued Atlantic overturning circulation even under climate extremes. \emph{Nature}. 
\bibitem{amoc2025b} (2025). Weakened Atlantic Meridional Overturning Circulation causes the historical North Atlantic Warming Hole. \emph{Communications Earth \& Environment}. 
\bibitem{volga2025} (2025). Modeling the Interannual and Seasonal Variability of the Caspian Sea Level Using ERA5 Atmospheric Data (1940–2024). \emph{Pure and Applied Geophysics}. 
\bibitem{caspian2025} (2025). World’s largest lake is drying up: Caspian Sea water levels drop to record low. 
\bibitem{ivt2025} Caspian Sea Level as an Indicator of Large-Scale Moisture Cycling in the Ocean-Atmosphere-Land System. 
\bibitem{nao2025} Past and future impact of North Atlantic teleconnection patterns on the hydroclimate of the Caspian catchment area. 
\bibitem{ml2025} Predictive Modeling of Global Warming Trends Using Machine Learning. 
\bibitem{shumakov2025} Roshydromet: the Caspian Sea shallowing will continue (2025). 
\bibitem{astana2026} Caspian Sea’s Decline Hits Economy, Infrastructure and Health (2026). 
\end{thebibliography}

\end{document}